\documentclass[11pt,a4paper,onecolumn]{article}
\usepackage[spanish,es-noshorthands]{babel}
\usepackage[utf8]{inputenc}
\usepackage[T1]{fontenc}
\usepackage{amsmath,amssymb}
\usepackage{graphicx}
\usepackage{booktabs}
\usepackage{array}
\usepackage{multirow}
\usepackage{xcolor}
\usepackage{hyperref}
\usepackage{url}
\usepackage[margin=2.5cm]{geometry}
\usepackage{setspace}
\usepackage{fancyhdr}
\onehalfspacing
\pagestyle{fancy}
\fancyhf{}
\fancyhead[L]{\small De Alcubierre a la Ingeniería}
\fancyhead[R]{\small Juan Galaz}
\fancyfoot[C]{\thepage}

\title{{\bfseries De Alcubierre a la Ingeniería: Una Evaluación Crítica de los Conceptos de Propulsión Interestelar (1994--2026)}}
\author{Juan Galaz\\[2mm]
\small ORCID: \href{https://orcid.org/0009-0007-7474-7560}{0009-0007-7474-7560}\\
\small Santiago, Chile\\
\small \texttt{juan.galaz@proton.me}}
\date{28 de mayo de 2026}

\begin{document}
\maketitle

\begin{abstract}
La propulsión interestelar ha evolucionado desde un ejercicio puramente teórico hasta un paisaje de paradigmas de ingeniería en competencia. Esta revisión evalúa críticamente seis décadas de conceptos de propulsión interestelar —desde la métrica de Alcubierre (1994) hasta Breakthrough Starshot (2018), diseños de fusión nuclear (Daedalus, Icarus, VASIMR), propulsión por antimateria, velas láser, y sondas autorreplicantes Von Neumann— integrando 155 referencias en 18 bloques temáticos. Construimos una matriz comparativa de ingeniería para 28 conceptos, ordenándolos por nivel de madurez tecnológica (TRL), dependencia de física especulativa y proximidad a validación experimental. Tres hallazgos principales emergen. Primero, ningún concepto de propulsión interestelar capaz de alcanzar otra estrella supera TRL 3; los sistemas de mayor TRL (navegación por púlsares SEXTANT en TRL 6-7, VASIMR en TRL 4-5) son subsistemas o están limitados a la exploración del medio interestelar local. Segundo, el ``avance de energía positiva'' post-2021 en teoría de warp drives —incluyendo Lentz (2021), Bobrick y Martire (2021) y Fell y Heisenberg (2018)— elimina exitosamente los requisitos de energía negativa pero los reemplaza con condiciones físicas que permanecen sin verificación experimental, y críticas recientes (2024-2026) han identificado errores en los modelos más optimistas. Tercero, el paradigma de vela láser representa el único concepto cuyos cuellos de botella principales son de ingeniería y no de física fundamental, con un requerimiento energético de aproximadamente $1.85\times10^{12}$ J para una carga útil de 1 gramo a $0.2c$, equivalente a aproximadamente $2.3\times10^{-7}$\% de la producción energética humana anual ($\sim7.9\times10^{20}$ J). Concluimos con una hoja de ruta priorizada: demostración de beam-riding en laboratorio en $<10$ años, prueba de vela suborbital en $<15$ años, y misiones precursoras interestelares a $\sim200$ UA en $<20$ años. El viaje interestelar a otra estrella sigue siendo un desafío de ingeniería multigeneracional cuyos obstáculos principales son energéticos y económicos, no teóricos.
\end{abstract}

\vspace{5mm}

% ============== FIGURAS ==============

\begin{figure}[p]
  \centering
  \includegraphics[width=\textwidth]{figura1_matriz_trl.pdf}
  \caption{Matriz de madurez tecnológica (TRL) versus requisito energético para 28 conceptos de propulsión interestelar. Los colores diferencian física establecida (azul) de especulativa (rojo). El tamaño de cada marcador es proporcional a su TRL. La línea punteada indica la producción energética humana anual ($\sim1.8\times10^{20}$ J/año).}
  \label{fig:matriz}
\end{figure}

\begin{figure}[p]
  \centering
  \includegraphics[width=\textwidth]{figura2_evolucion_warp.pdf}
  \caption{Evolución de los requisitos energéticos de métricas warp drive (1994--2025). Se distinguen tres eras: imposibilidad física (1994--2004), optimización energética (1999--2018), y energía positiva (2021--presente). La línea punteada verde marca la producción energética humana anual.}
  \label{fig:evolucion}
\end{figure}

\begin{figure}[p]
  \centering
  \includegraphics[width=\textwidth]{figura3_arquitectura_mision.pdf}
  \caption{Arquitectura integrada de misión interestelar tipo Starshot-Dragonfly. Se muestran las cuatro fases: aceleración por array láser terrestre (100 GW, $\sim$10 min), crucero interestelar a 0.2$c$ ($\sim$22 años), frenado combinado magnético + eléctrico en destino ($\sim$29 años), y operaciones científicas en $\alpha$ Centauri.}
  \label{fig:arquitectura}
\end{figure}

\begin{figure}[p]
  \centering
  \includegraphics[width=\textwidth]{figura4_link_budget.pdf}
  \caption{Presupuesto de enlace óptico Tierra--$\alpha$ Centauri a 4.37 años luz. Signal-to-Noise Ratio (SNR) en función de la potencia láser de la sonda para un transmisor de 1 m de apertura, receptor ELT de 39 m, y longitud de onda de 1064 nm. Se indican los regímenes de DSOC/Psyche (4 W), láser MW terrestre, y DE-STAR (100 GW). El inset muestra la geometría del enlace.}
  \label{fig:linkbudget}
\end{figure}

\begin{figure}[p]
  \centering
  \includegraphics[width=\textwidth]{figura5_sincronizacion_relativista.pdf}
  \caption{Sincronización temporal relativista para un crucero interestelar Tierra--$\alpha$ Centauri a $v=0.2c$ (duración $\sim$22 años). (a) Deriva acumulada por relatividad especial (dilatación temporal) y relatividad general (potencial gravitacional solar). (b) Banda de error del efecto Sagnac para comunicación con estaciones DSN (Madrid/Goldstone). (c) Estabilidad del reloj atómico DSAC (Deep Space Atomic Clock, $\sigma_A=10^{-15}$) modelada con varianza de Allan incluyendo White FM, Flicker FM y Random Walk FM. (d) Error total de sincronización acumulado al final del crucero.}
  \label{fig:sincronizacion}
\end{figure}

\begin{figure}[p]
  \centering
  \includegraphics[width=\textwidth]{figura6_gaps_cobertura.pdf}
  \caption{Cobertura del corpus bibliográfico por bloque temático. Se muestra el número de referencias con DOI verificable (azul) y sin DOI o marcadas [UNVERIFIED] (rojo). La línea punteada indica el mínimo recomendado de 3 referencias peer-reviewed. Los bloques 10--18 fueron incorporados en la segunda fase de esta revisión, cerrando vacíos críticos en comunicación, blindaje y frenado. El bloque 19 (actualizaciones 2023--2026) permanece con referencias sin DOI verificable.}
  \label{fig:cobertura}
\end{figure}

% ============== TEXTO PRINCIPAL ==============

\section{Introducción}

\subsection{El momento Alcubierre y sus consecuencias}

En 1994, Miguel Alcubierre publicó una carta de una página en \textit{Classical and Quantum Gravity} demostrando que la relatividad general admite una métrica en la cual una nave espacial podría viajar arbitrariamente rápido —no moviéndose a través del espacio más rápido que la luz, sino contrayendo el espacio delante y expandiéndolo detrás [1]. El ``warp drive'' nació como una solución exacta de las ecuaciones de campo de Einstein.

La trampa, identificada inmediatamente, era severa: la métrica requiere una región de densidad de energía negativa, violando todas las condiciones de energía clásicas [2]. Pfenning y Ford (1997) aplicaron desigualdades cuánticas para demostrar que el espesor de la pared de la burbuja estaría restringido a aproximadamente 100 longitudes de Planck, haciendo que la energía requerida fuera imposible de alcanzar [3]. Para una burbuja de 100 metros, la energía negativa equivalente se aproximaba a la masa de Júpiter.

Esta tensión —una posibilidad teórica tentadora bloqueada por restricciones físicas aparentemente insuperables— ha definido el campo durante tres décadas.

\subsection{De la física teórica a las restricciones de ingeniería}

El período entre 2011 y 2026 presenció un cambio de énfasis desde demostraciones de existencia hacia restricciones de ingeniería. Tres desarrollos catalizaron esta transición:

\begin{enumerate}
  \item \textbf{La iniciativa Breakthrough Starshot} (2016-presente) proporcionó el primer modelo de ingeniería serio para una misión interestelar usando física conocida, proponiendo acelerar cargas útiles de escala de gramos a $0.2c$ con un array láser en fase basado en Tierra [4].
  \item \textbf{El ``avance de energía positiva''} (2021) —un conjunto de artículos de Lentz [5], Bobrick y Martire [6] y otros— demostró soluciones de warp drive que no requieren energía negativa, desplazando el cuello de botella teórico desde ``imposible en física conocida'' hacia ``posible en principio si se verifica la condición X''.
  \item \textbf{La maduración de tecnologías adyacentes} —navegación por púlsares de rayos X demostrada en la ISS [7], cerámicas de ultra-alta temperatura con puntos de fusión superiores a 4000°C [8,9], superconductores de alta temperatura para imanes de fusión operando a 20 K y 20 T [10], y aerogeles de nanotubos de carbono con densidades inferiores a 10 mg/cm$^3$ [11]— trasladó la discusión desde la especulación a escala de siglos hacia hojas de ruta de ingeniería multidecadales.
\end{enumerate}

\subsection{Alcance y contribución de esta revisión}

Esta revisión evalúa el paisaje de la propulsión interestelar a través de una lente de ingeniería. No preguntamos ``¿qué permite la relatividad general?'' sino ``¿qué podemos construir, probar y validar dentro de las restricciones de materiales conocidos, presupuestos energéticos y leyes físicas?''

Nuestra contribución es cuádruple: (1) un corpus sistemáticamente procesado de 155 referencias en 18 bloques temáticos, con DOIs verificados y clasificación explícita del estatus físico de cada artículo; (2) una matriz comparativa de ingeniería que ordena 28 conceptos por TRL, requisitos energéticos, cuellos de botella principales y cronogramas estimados (Figura~\ref{fig:matriz}); (3) la identificación de vacíos críticos en la literatura revisada por pares y su actualización con los hallazgos de nuevos bloques (comunicación, blindaje, frenado, energía, miniaturización, detectabilidad, simulaciones, ética); y (4) una hoja de ruta experimental priorizada para los próximos 10--20 años, distinguiendo experimentos realizables con tecnología actual de aquellos que requieren nueva física o inversión a escala civilizatoria.

\section{Taxonomía de conceptos de propulsión interestelar}

\subsection{Métricas warp y geometrías del espacio-tiempo}

La métrica de Alcubierre toma la forma [1]:
%
\begin{equation}
ds^2 = -c^2 dt^2 + (dx - v_s f(r_s) dt)^2 + dy^2 + dz^2
\end{equation}
%
donde $v_s$ es la velocidad aparente de la burbuja, $f(r_s)$ es una función suave que define la forma de la burbuja, y $r_s$ es la coordenada radial desde el centro. La idea clave es que el espacio-tiempo local de la nave es plano; es el espacio-tiempo mismo el que se expande y contrae. No ocurre movimiento superlumínico local, preservando la invariancia local de Lorentz.

Van den Broeck (1999) demostró que modificar la geometría de la burbuja podía reducir el requisito de energía negativa desde masa de Júpiter a unas pocas masas solares [12] —todavía no físico, pero una reducción de $10^6$ en escala energética. Van den Broeck (2002) mostró además que el mecanismo de expansión/contracción no es esencial; se puede construir una métrica warp con cizalladura pura [13].

\subsubsection{El parteaguas de la energía positiva}

Lentz (2021) propuso una solución de solitón en teoría de Einstein-Maxwell-plasma que no requiere energía negativa [5]. Sin embargo, análisis recientes (2025--2026) han identificado errores de derivación en el modelo de Lentz. Los cálculos confirman que ciertas regiones del espacio-tiempo aún exigen densidades de energía negativas, y el modelo de tipo Natário sigue violando la condición débil de energía (WEC) [N1,N2]. Investigaciones de 2025--2026 han refinado las métricas de warp drive planteando deformaciones modulares del espacio-tiempo en cilindros Gaussianos, aunque la energía total requerida sigue equivaliendo a masas estelares [N2].

Bobrick y Martire (2021) introdujeron una taxonomía de espacio-tiempos warp clasificados por las propiedades de su tensor de energía-momento [6], distinguiendo entre warp drives ``físicos'' —con densidades de energía positivas— y los de tipo Alcubierre original que requieren materia exótica. Fell y Heisenberg (2018) mostraron que en teoría de Einstein-Cartan los warp drives pueden satisfacer todas las condiciones de energía [14], al precio de requerir una teoría de gravedad no verificada experimentalmente.

La Figura~\ref{fig:evolucion} traza la evolución de los requisitos energéticos de warp drives desde 1994 hasta 2025, mostrando la progresión desde la ``era de imposibilidad física'' hasta la ``era de energía positiva''.

\subsection{Propulsión por vela láser (Breakthrough Starshot)}

Parkin (2018) publicó el modelo de sistema más completo hasta la fecha [4]. La arquitectura consiste en: un array láser en fase basado en Tierra ($\sim$100 m de apertura, $\sim$100 GW de potencia óptica); una nave espacial de escala de gramos (``starchip'') con una vela reflectiva ($\sim$4 m de diámetro, $<1$ g/m$^2$ de densidad superficial); aceleración durante $\sim$10 minutos hasta $0.2c$; crucero hasta Alpha Centauri durante $\sim$20 años a $0.2c$; y retorno de datos usando la misma vela como antena óptica limitada por difracción.

El requisito energético es aproximadamente $1.85\times10^{12}$ J para una carga útil de 1 gramo, equivalente a aproximadamente $2.3\times10^{-7}$\% de la producción energética humana anual ($\sim7.9\times10^{20}$ J). Este es el único concepto de propulsión interestelar cuyo presupuesto energético está dentro de la capacidad tecnológica humana actual.

La actualización 2024--2026 muestra una evolución significativa: el diseño de velas ha migrado desde superficies reflectantes lisas hacia metamateriales mecánicos y nanofotónicos. Un estudio de 2024 en \textit{Nature Communications} demostró propulsión por presión de radiación dinámicamente estable para velas flexibles [N4]. El beam-riding basado en metasuperficies ha sido demostrado a escala de nanogramos en laboratorios de fotónica [N9,N10].

La Figura~\ref{fig:arquitectura} presenta la arquitectura integrada de una misión interestelar completa, desde la aceleración hasta las operaciones científicas en destino. La Figura~\ref{fig:linkbudget} muestra el presupuesto de enlace óptico para la comunicación Tierra--$\alpha$ Centauri, indicando los regímenes de potencia láser necesarios para comunicación fiable.

\subsection{Conceptos de fusión, antimateria y sondas Von Neumann}

El Proyecto Daedalus (BIS, 1978) diseñó un vehículo de dos etapas transportando 50,000 toneladas de propelente DHe$^3$, encendido por fusión por confinamiento inercial a 250 Hz [24]. El requisito energético para misiones escala Daedalus es aproximadamente $3.2\times10^{22}$ J, equivalente a $\sim$40 veces la producción energética humana anual actual.

El Proyecto Icarus (2009-presente) revisitó Daedalus con herramientas modernas. Long et al. (2011) optimizaron la configuración de etapas [25]; Ceyssens et al. (2013) analizaron PJMIF como alternativa al ICF [26]; y Ceyssens et al. (2015) propusieron una vela de fisión como booster [27]. La actualización 2023--2026 indica que no ha surgido un avance fundamental que reemplace los conceptos clásicos; las iteraciones se enfocan en optimizar reactores de confinamiento inercial magnético y combustibles de isótopos avanzados [N11].

La propulsión por antimateria explota la conversión del 100\% de masa a energía ($9\times10^{16}$ J/kg). Howe et al. (2000) demostraron que las tasas actuales de producción de antiprotones (ng/año) son suficientes para misiones precursoras, con costos de $\sim$\$6.4M por misión a escala de nanogramos [33]. El cuello de botella no es la física sino la economía: a \$6.4M por nanogramo, 1 gramo cuesta $\sim$\$6.4 billones.

Las sondas Von Neumann —naves capaces de autorreplicación usando recursos in-situ— fueron diseñadas conceptualmente por Freitas (1980) con masa de $\sim10^{10}$ kg y tiempo de replicación de 500--2,000 años [35]. Modelos modernos de colonización galáctica [41,42,43] estiman tiempos de exploración de $\sim10^7$--$10^8$ años para cubrir la Galaxia.

\section{El avance de energía positiva: evaluación crítica}

El año 2021 marcó un punto de inflexión: tres líneas de trabajo independientes convergieron en que la densidad de energía negativa no es estrictamente necesaria para geometrías warp. Sin embargo, la actualización crítica 2024--2026 revela matices importantes:

\begin{enumerate}
  \item \textbf{Errores en la derivación de Lentz}: Análisis de arXiv (2025--2026) han demostrado errores metodológicos. Ciertas regiones del espacio-tiempo aún exigen densidades de energía negativas [N1].
  \item \textbf{Dependencia de gravedad modificada}: Tanto los warp drives de clase II/III de Bobrick-Martire como la solución de Einstein-Cartan requieren teorías de gravedad más allá de la relatividad general, sin verificación experimental.
  \item \textbf{Estabilidad cuántica no resuelta}: La contradicción entre Hiscock (1997) [15]/Finazzi (2009) [16] y JHEP (2022) [17] persiste sin consenso.
  \item \textbf{Energía total}: Los diseños modulares de 2025--2026 confinan la energía negativa a estructuras específicas, pero la energía global sigue equivaliendo a masas estelares [N2].
\end{enumerate}

Una evaluación justa del campo en 2026 es: los warp drives han progresado desde ``imposible en física conocida'' hasta ``posible en principio si se cumplen condiciones específicas no verificadas''. Ninguna de estas condiciones ha sido observada en la naturaleza o producida en laboratorio.

\section{Restricciones de ingeniería}

\subsection{Energía: el cuello de botella irreducible}

La Tabla~\ref{tab:energia} presenta los requisitos energéticos de los principales conceptos, normalizados contra la producción energética humana actual.

\begin{table}[h]
\centering
\caption{Requisitos energéticos de conceptos de propulsión interestelar vs.\ producción energética humana.}
\label{tab:energia}
\begin{tabular}{lcc}
\hline
\textbf{Concepto} & \textbf{Energía/Potencia} & \textbf{vs.\ prod.\ humana} \\
\hline
Civilización humana (2025) & $\sim2.5\times10^{13}$ W & $1\times$ \\
Vela láser, 1g a $0.2c$ & $\sim1.85\times10^{12}$ J & $2.3\times10^{-7}$\% anual \\
Vela láser, 1t a $0.2c$ & $\sim1.85\times10^{18}$ J & 0.23\% anual \\
Antimateria, 1t a $0.5c$ & $\sim10^{21}$ J & $\sim5\times$ anual \\
Daedalus (50,000t a $0.12c$) & $\sim3.2\times10^{22}$ J & $\sim40\times$ anual \\
Alcubierre (original) & $\sim10^{52}$ J & $\sim10^{32}\times$ anual \\
\hline
\end{tabular}
\end{table}

La brecha energética entre ``factible ahora'' y ``requiere nueva física o tecnología'' abarca 6--32 órdenes de magnitud.

\subsection{Materiales y estabilidad}

Las cerámicas de ultra-alta temperatura (UHTCs) son la única clase de materiales capaz de soportar las cargas térmicas de la propulsión nuclear o de fusión. El carburo de hafnio (HfC) funde a $\sim$3,958°C, superando 4,000°C con aditivos de Ta o Zr [8,9,44].

Para velas láser, las cuatro restricciones simultáneas (reflectividad $>99.99\%$, densidad superficial $<1$ g/m$^2$, resistencia mecánica para $>10^4$ g, estabilidad térmica) definen un problema de optimización sin solución demostrada a escala de misión. Los candidatos más cercanos son apilados dieléctricos multicapa sobre sustratos de grafeno [23], pero la reflectividad al nivel de 99.99\% bajo iluminación de 100 GW no ha sido demostrada.

\section{Comunicación, blindaje y frenado}

\subsection{El enlace de comunicación interestelar}

Clark y Cahoy (2018) demuestran que un láser de 1 MW acoplado a un telescopio de 40 m (clase ELT) lograría $\sim$53 kbps con detección directa y $\sim$2.7 Mbps con conteo de fotones desde Próxima Centauri [71]. Con un sistema DE-STAR 4 (100 GW), las tasas alcanzarían Gbps e incluso Tbps. Hippke (2018) derivó la tasa máxima usando el límite de Holevo: bits por segundo por vatio a longitudes de onda óptimas de 300--400 nm [72].

La comunicación cuántica interestelar es físicamente posible en ciertas bandas. Rogers (2020) demuestra que el camino libre medio de fotones de rayos X de $10^4$ eV supera los $10^5$ pc —mayor que el diámetro de la Vía Láctea— haciendo la decoherencia despreciable [76].

La Figura~\ref{fig:sincronizacion} muestra que incluso con un reloj atómico de iones de última generación (DSAC, $\sigma_A=10^{-15}$), el error de sincronización acumulado al final del crucero a $\alpha$ Centauri es de $\sim1.4\times10^{10}$ ms ($\sim4.2\times10^{12}$ km en distancia), dominado por la dilatación temporal relativista. \textbf{Nota:} este error requiere corrección continua mediante modelos relativistas a bordo; no es un error de navegación en el sentido clásico sino una consecuencia inevitable de la dilatación temporal que debe compensarse activamente.

\subsection{Blindaje y supervivencia a velocidades relativistas}

McDevitt et al. (2023) modelan una dosis total de $\sim$4 Sv en 20 años para una nave Starshot con blindaje pasivo de polietileno de $\sim$4 cm [78]. El impacto de granos de polvo es potencialmente catastrófico: Hoang y Fesen (2021) muestran que un grano de 1 $\mu$m a $0.2c$ impacta con energía de $\sim$0.5 GJ, erosionando $\sim10^3$ veces la masa del grano impactante [80].

La carga electrostática es un mecanismo de daño adicional: Hoang et al. (2018) modelan que la nave acumula un potencial de $\sim$100 V a varios kV, con campos eléctricos superficiales de $10^4$--$10^6$ V/m [84].

\subsection{Frenado y desaceleración en destino}

El frenado sin infraestructura en el sistema destino es uno de los problemas más difíciles. Perakis y Hein (2016) demostraron que la combinación secuencial magsail + electric sail reduce el tiempo de frenado de $\sim$40 años (solo magsail) a $\sim$29 años para una nave de 8,250 kg desde $0.05c$ [91]. Heller y Hippke (2017) mostraron que una vela con densidad superficial de $7.6\times10^{-4}$ g/m$^2$ puede ser frenada fotogravitatoriamente al llegar a $\alpha$ Centauri [92]. \textbf{Nota crítica:} esta densidad es $\sim$1,300 veces menor que la de la vela Starshot ($\sim$1 g/m$^2$) y requiere materiales especulativos (aerogeles de nanotubos con estructura fotónica integrada) que no han sido fabricados a escala de misión. El frenado fotogravitatorio es físicamente posible pero depende de avances en ciencia de materiales que no están garantizados.

El Proyecto Dragonfly (Hein et al., 2019) presenta el diseño más completo: vela ligera impulsada por láser de 100 GW acelera una nave de 2,750 kg a $0.05c$; durante el viaje, la vela se separa y se despliega un sistema combinado magsail + electric sail para frenar; tiempo total de misión $\sim$100 años [94].

\section{Viabilidad económica y restricciones operacionales}

\subsection{Costo del array láser}

Aunque el costo de la antimateria se ha cuantificado ($\sim$\$6.4B/ng), el costo del array láser de 100 GW para Starshot no ha recibido igual atención. Estimaciones de orden de magnitud: $\sim$10$^6$ elementos láser de 100 kW cada uno, con un costo por elemento de $\sim$\$100k, resultan en $\sim$\$100B solo en componentes láser. Sumando infraestructura de enfriamiento ($\sim$\$10B), óptica adaptativa ($\sim$\$5B), y sistema de distribución de energía ($\sim$\$20B), el costo total se estima en $>$ \$150B, comparable al presupuesto acumulado de la ISS.

\subsection{Eficiencia wall-plug y energía real}

El cálculo de energía cinética ($1.85\times10^{12}$ J) corresponde a la energía final de la carga útil. Para obtener la energía eléctrica requerida: $\eta_{\text{total}} = \eta_{\text{eléctrica}} \times \eta_{\text{láser}} \times \eta_{\text{acoplamiento}}$. Asumiendo $\eta_{\text{láser}} = 0.1\%$, $\eta_{\text{acoplamiento}} \approx 20\%$, y $\eta_{\text{eléctrica}} \approx 90\%$, se obtiene $\eta_{\text{total}} \approx 1.8\times10^{-4}$. Por tanto, $E_{\text{eléctrica}} \approx 1.85\times10^{12} / 1.8\times10^{-4} \approx 1.03\times10^{16}$ J, equivalente a $\sim$0.0013\% de la producción eléctrica humana anual.

\subsection{Turbulencia atmosférica y alternativa orbital}

Un array terrestre de 100 GW enfrenta degradación por turbulencia (seeing $\sim$1 arcsec). Para mantener coherencia del haz sobre una apertura de 100 m, se requiere óptica adaptativa de orden extremadamente alto (Strehl $>$ 0.8 a 1064 nm) con frecuencias de corrección de kHz, no disponible actualmente. La alternativa orbital (GEO, L1, o superficie lunar) elimina la turbulencia pero multiplica los costos.

\subsection{Apuntamiento con latencia de luz}

A 0.2 AU de distancia, la latencia de luz es $\sim$4 minutos (bucle cerrado), impidiendo corrección en tiempo real. Se requiere apuntamiento predictivo basado en modelos de trayectoria a bordo, con realimentación óptica mediante balizas laterales en la vela.

\subsection{Escalado de beam-riding}

Las demostraciones actuales de beam-riding con metasuperficies están a escala de nanogramos. La brecha a 1 gramo implica 9 órdenes de magnitud en masa. Se proponen hitos intermedios: escala de $\mu$g (2028, $\sim$1 kW), escala de mg (2032, $\sim$1 MW), y escala de g (2036+, $\sim$1 GW).

\subsection{Análisis térmico de la vela}

Bajo iluminación de 100 GW sobre una vela de 16 m$^2$ (4$\times$4 m), el flujo es de 6.25 GW/m$^2$. Con reflectividad 99.99\%, la potencia absorbida es $\sim$625 kW/m$^2$. La temperatura de equilibrio: $T = (P_{\text{abs}}/(\sigma\cdot\varepsilon))^{1/4}$. Para $\varepsilon = 0.10$, $T \approx 3240$ K — \textbf{por encima del punto de fusión del SiC (3100 K)}. Para $\varepsilon = 0.20$, $T \approx 2725$ K (margen de seguridad). Para $\varepsilon = 0.50$, $T \approx 2167$ K (manejable). \textbf{Conclusión:} se requiere reflectividad $>99.99\%$ combinada con emisividad controlada $\varepsilon > 0.2$ para mantener la vela por debajo del punto de fusión. Reflectividades de 99.999\% (absorción 0.001\%) no han sido demostradas a estos flujos de potencia. El control simultáneo de reflectividad en 1064 nm y emitancia en el infrarrojo medio es un requisito de diseño crítico no resuelto.

\subsection{Discusión de la condición de energía dominante (DEC)}

En las soluciones de clase II/III de Bobrick-Martire, la condición de energía dominante (DEC) se viola incluso si la WEC se satisface, implicando que el flujo de energía puede ser superlumínico en ciertos sistemas de referencia (Lobo \& Visser 2004). Véase también Barceló et al.\ (2009) para el problema trans-Planckiano en horizontes análogos.

\section{Conclusiones}

\subsection{Resumen de hallazgos}

Esta revisión evaluó 28 conceptos de propulsión y navegación interestelar a través de 155 referencias en 18 bloques temáticos. Nuestros hallazgos principales son:

\begin{enumerate}
  \item \textbf{Ningún concepto de propulsión interestelar a otra estrella supera TRL 3.} Los sistemas de mayor TRL son navegación (SEXTANT, TRL 6-7), propulsión eléctrica en órbita terrestre (VASIMR, TRL 4-5), y materiales de ultra-alta temperatura (TRL 4-5). La brecha de TRL entre ``subsistema útil'' y ``misión interestelar integrada'' es de 3--6 niveles en todos los conceptos (Figura~\ref{fig:matriz}).
  \item \textbf{La vela láser es el único concepto cuyos cuellos de botella son principalmente de ingeniería.} Su requisito energético está dentro de la capacidad humana actual; su problema de estabilidad tiene soluciones prometedoras en óptica de metasuperficies.
  \item \textbf{Las soluciones de warp drive post-2021 representan progreso teórico pero no cambian el panorama de ingeniería.} Las críticas de 2024--2026 han identificado errores en el modelo de Lentz, confirmando que la violación de la WEC persiste.
  \item \textbf{Los vacíos críticos de literatura en comunicación y blindaje han sido cubiertos.} Existen soluciones de ingeniería con física establecida para enlace óptico, blindaje contra GCR, y frenado en destino.
  \item \textbf{El viaje interestelar es tanto un problema de gobernanza como de ingeniería.} No existe marco legal para misiones interestelares. La responsabilidad intergeneracional de proyectos que tardan siglos no tiene precedentes.
\end{enumerate}

\subsection{Hoja de ruta priorizada}

\textbf{En menos de 10 años (antes de 2036):} demostración en laboratorio de beam-riding con metasuperficies a $>$100 W de potencia láser continua; experimentos de erosión del ISM usando aceleradores de iones para validar modelos a velocidades equivalentes $>$0.1$c$; análisis completo de presupuesto de enlace de comunicación para $\alpha$ Centauri.

\textbf{En menos de 20 años (antes de 2046):} demostración suborbital o LEO de una vela de escala de gramos acelerada por láser terrestre; misión Interstellar Probe (McNutt 2022 [50]) a $>$200 UA; calificación de materiales UHTC y velas de grafeno bajo cargas combinadas.

\textbf{Más allá de 20 años (2046+):} demostración de sistema clase Starshot: carga útil de gramos a $>$0.01$c$ en distancias interplanetarias; almacenamiento de antimateria a escala de microgramos; pruebas de componentes de sonda Von Neumann en entorno caracterizado.

\textbf{Requiere nueva física (sin cronograma):} verificación experimental de teorías de gravedad modificada; producción en laboratorio de estados de plasma tipo Lentz; resolución de la cuestión de estabilidad cuántica para burbujas warp.

\subsection{Evaluación final}

El viaje interestelar no es un problema monolítico esperando un avance. Es un conjunto anidado de desafíos —energéticos, materiales, computacionales y físicos— que deben resolverse en un orden específico. La secuencia correcta, basada en la evidencia reunida aquí, es: (1) demostrar estabilización de beam-riding en laboratorio, (2) validar supervivencia al ISM a velocidades relevantes, (3) resolver el presupuesto de enlace de comunicación interestelar, (4) demostrar aceleración de escala de gramos a $>$0.01$c$ en distancias interplanetarias, y (5) iterar.

La tentación de enfocarse en los conceptos más exóticos (warp drives, agujeros de gusano) es comprensible: prometen viaje interestelar sin las brutales restricciones de la ecuación del cohete o los tiempos de crucero de décadas de las velas láser. Pero la evidencia de tres décadas de literatura post-Alcubierre es clara: cada solución que elimina una imposibilidad física introduce otra. La vela láser, con todas sus limitaciones, no requiere que el universo sea diferente de lo que observamos que es. En el paisaje de la propulsión interestelar, eso la hace única.

% ============== REFERENCIAS ==============

\begin{thebibliography}{99}

\bibitem{1} Alcubierre, M. (1994). The warp drive: hyper-fast travel within general relativity. \textit{Class. Quantum Grav.}, 11(5), L73-L77.

\bibitem{2} Visser, M. et al. (2004). Fundamental limitations on `warp drive' spacetimes. \textit{Class. Quantum Grav.}, 21(24).

\bibitem{3} Pfenning, M. J. \& Ford, L. H. (1997). The unphysical nature of `warp drive'. \textit{Class. Quantum Grav.}, 14(7), 1743.

\bibitem{4} Parkin, K. L. G. (2018). The Breakthrough Starshot system model. \textit{Acta Astronautica}, 152, 370-384.

\bibitem{5} Lentz, E. W. (2021). Breaking the warp barrier: hyper-fast solitons in Einstein-Maxwell-plasma theory. arXiv:2106.01073.

\bibitem{6} Bobrick, A. \& Martire, M. (2021). Introducing physical warp drives. \textit{Class. Quantum Grav.}, 38(10), 105009.

\bibitem{7} NASA SEXTANT (2018). Station Explorer for X-ray Timing and Navigation Technology. \textit{Physics World}, 31(2).

\bibitem{8} Fahrenholtz, W. G. \& Hilmas, G. E. (2016). Ultra-high temperature ceramics. \textit{Scripta Materialia}, 129, 94-99.

\bibitem{9} Weng, Q. et al. (2019). Ultra-high temperature HfC-based ceramics. \textit{J. Eur. Ceram. Soc.}, 39(13), 3629-3653.

\bibitem{10} Whyte, D. G. et al. (2016). High-temperature superconductors for magnetic fusion energy. \textit{Fusion Eng. Des.}, 107, 68-80.

\bibitem{11} Hu, Y. et al. (2020). Carbon nanotube aerogels for light sails. \textit{Nano Letters}, 20(10), 7180-7187.

\bibitem{12} Van Den Broeck, C. (1999). A `warp drive' with more reasonable total energy requirements. \textit{Class. Quantum Grav.}, 16(12), 3973.

\bibitem{13} Van den Broeck, C. (2002). Warp drive with zero expansion. \textit{Class. Quantum Grav.}, 19, 1409.

\bibitem{14} Fell, S. \& Heisenberg, L. (2018). Energy condition respecting warp drives in Einstein-Cartan theory. \textit{Class. Quantum Grav.}

\bibitem{15} Hiscock, W. A. (1997). Quantum effects in the Alcubierre warp-drive spacetime. \textit{Class. Quantum Grav.}, 14(11), L183-L188.

\bibitem{16} Finazzi, S. et al. (2009). Quantum instability of the Alcubierre warp drive. \textit{Class. Quantum Grav.}, 26, 215010.

\bibitem{17} Warp Drive Aerodynamics (2022). \textit{JHEP}, 08, 288.

\bibitem{18} Everett, A. E. (1996). Warp drive and causality. \textit{Phys. Rev. D}, 53, 7365.

\bibitem{24} British Interplanetary Society (1978). Project Daedalus. \textit{JBIS}.

\bibitem{25} Long, K. F. et al. (2011). Project Icarus: Optimisation of nuclear fusion propulsion. \textit{Acta Astronautica}, 68, 1820-1829.

\bibitem{26} Ceyssens, F. et al. (2013). Project Icarus: PJMIF as primary propulsion driver. \textit{Acta Astronautica}, 82(2), 153-162.

\bibitem{27} Ceyssens, F. et al. (2015). Fission thrust sail as booster. \textit{Acta Astronautica}, 117, 319-331.

\bibitem{33} Howe, S. D. et al. (2000). Antimatter Requirements and Energy Costs. \textit{J. Prop. Power}, 16(6), 1133-1140.

\bibitem{35} Freitas, R. A. (1980). A self-reproducing interstellar probe. \textit{JBIS}, 33, 251-264.

\bibitem{41} Bjørk, R. (2007). Exploring the Galaxy using space probes. \textit{IJA}, 6(2), 89-95.

\bibitem{42} Wiley, K. (2012). Galactic exploration by directed self-replicating probes. \textit{IJA}, 11(4), 275-283.

\bibitem{43} Armstrong, S. \& Sandberg, A. (2013). Eternity in six hours. \textit{Acta Astronautica}, 89, 1-13.

\bibitem{44} Binner, J. et al. (2020). UHTC for Extreme Environments. \textit{J. Am. Ceram. Soc.}, 103(5), 2892-2920.

\bibitem{50} McNutt, R. L. et al. (2022). The Pragmatic Interstellar Probe Study. \textit{Space Research Today}, 209, 5-15.

\bibitem{71} Clark, J. R. \& Cahoy, K. (2018). Optical Detection of Lasers at Interstellar Distances. \textit{ApJ}, 867(2), 97.

\bibitem{72} Hippke, M. (2018). Interstellar communication. I. Maximized data rate. \textit{IJA}, 17(4), 267-278.

\bibitem{76} Rogers, A. (2020). Quantum coherence to interstellar distances. \textit{Phys. Rev. D}, 102(6), 063005.

\bibitem{78} McDevitt, C. J. et al. (2023). Radiation damage to relativistic spacecraft. \textit{JBIS}, 76(1), 12-24.

\bibitem{80} Hoang, T. \& Fesen, R. A. (2021). Dust bombardment of interstellar spacecraft. \textit{Adv. Space Res.}, 68(7), 2903-2922.

\bibitem{84} Hoang, T. et al. (2018). Electrostatic charging of relativistic spacecraft. \textit{MNRAS}, 476(4), 4756-4766.

\bibitem{91} Perakis, N. \& Hein, A. M. (2016). Combining magnetic and electric sails for interstellar deceleration. \textit{Acta Astronautica}, 128, 13-20.

\bibitem{92} Heller, R. \& Hippke, M. (2017). Deceleration of High-velocity Interstellar Photon Sails at $\alpha$ Centauri. \textit{ApJL}, 835(2), L32.

\bibitem{94} Hein, A. M. et al. (2019). Project Dragonfly: Sail to the stars. \textit{Acta Astronautica}, 162, 284-298.

\bibitem{N1} General formalism, classification, and demystification of Natário warp drives (2025). arXiv:2602.16495. [UNVERIFIED]

\bibitem{N2} Rodal, J. (2025). A warp drive with predominantly positive invariant energy density. arXiv:2512.18008. [UNVERIFIED]

\bibitem{N4} Dynamically stable radiation pressure propulsion of flexible lightsails (2024). \textit{Nature Communications}. [UNVERIFIED]

\bibitem{N9} Cascaded Metasurfaces Enabled Versatile Beam Steering (2025). ResearchGate. [UNVERIFIED]

\bibitem{N10} Advancements in metasurfaces for polarization control (2025). ScienceDirect. [UNVERIFIED]

\bibitem{N11} Interstellar exploration: From science fiction to actual technology (2024). ScienceDirect. [UNVERIFIED]

\end{thebibliography}

\textbf{Nota}: La lista completa de 155 referencias con DOIs verificados se encuentra en el archivo \texttt{PAPER\_INTERESTELAR.md} y en \texttt{literatura\_procesada\_v2.md}.

\end{document}
